\documentclass[12pt,a4paper]{exam}
\usepackage[utf8]{inputenc}
\usepackage{amsmath,amssymb,amsfonts}
\usepackage{geometry}
\usepackage{enumitem}
\usepackage{graphicx}
\usepackage{hyperref}
\usepackage{xcolor}
\geometry{a4paper, top=2cm, bottom=2cm, left=2.5cm, right=2.5cm}
\hypersetup{colorlinks=true, linkcolor=blue, urlcolor=blue}
\renewcommand{\thequestion}{\textbf{\arabic{question}}}
\renewcommand{\questionlabel}{\thequestion.}
\pointname{ marks}
\pointsdroppedatright
\bracketedpoints
\setlength{\parindent}{0pt}
\setlength{\emergencystretch}{3em}
\allowdisplaybreaks
\newsavebox{\schmathbox}
\newcommand{\fitmath}[1]{%
  \sbox{\schmathbox}{$\displaystyle #1$}%
  \ifdim\wd\schmathbox>\linewidth
    \resizebox{\linewidth}{!}{\usebox{\schmathbox}}%
  \else
    \usebox{\schmathbox}%
  \fi}

\begin{document}

\thispagestyle{empty}
\begin{center}
  \textbf{\LARGE Diag Solutions} \\[0.3cm]
  \rule{\textwidth}{0.5pt}
\end{center}

\vspace{0.3cm}

\begin{questions}
\question \textbf{1.} Two concentric circles are of radii $5$ cm and $3$ cm. The length of the chord of the larger circle which touches the smaller circle is:

\textbf{Answer:}
(B)

\textbf{Solution:}
\begin{enumerate}
  \item Let $O$ be the center, $AB$ be the chord of the larger circle touching the smaller circle at $M$.
  \item In $\triangle OMA$, $\angle OMA = 90^\circ$, $OA = 5$ cm, $OM = 3$ cm.
  \item By Pythagoras: $AM^2 = OA^2 - OM^2 = 25 - 9 = 16$. So $AM = 4$ cm.
  \item The perpendicular from the center bisects the chord, so $AB = 2 \times AM = 8$ cm.
\end{enumerate}
\textbf{Derivation:}
\begin{center}
\fitmath{AB = 2\sqrt{5^2 - 3^2} = 2\sqrt{16} = 8}
\end{center}
\hfill \textit{Circles — Concentric circles}
\vspace{0.3cm}

\question \textbf{2.} If $TP$ and $TQ$ are two tangents to a circle with center $O$ so that $\angle POQ = 110^\circ$, then $\angle PTQ$ is equal to:

\textbf{Answer:}
(C)

\textbf{Solution:}
\begin{enumerate}
  \item In the quadrilateral $\text{OPTQ}$, the sum of angles is $360^\circ$.
  \item Tangents are perpendicular to radii at contact points, so $\angle OPT = 90^\circ$ and $\angle OQT = 90^\circ$.
  \item $\angle PTQ + \angle POQ + \angle OPT + \angle OQT = 360^\circ$.
  \item $\angle PTQ + 110^\circ + 90^\circ + 90^\circ = 360^\circ \implies \angle PTQ = 360^\circ - 290^\circ = 70^\circ$.
\end{enumerate}
\textbf{Derivation:}
\begin{center}
\fitmath{\angle PTQ = 180^\circ - 110^\circ = 70^\circ}
\end{center}
\hfill \textit{Circles — Tangents from an external point}
\vspace{0.3cm}

\question \textbf{3.} Assertion (A): The length of a tangent drawn from an external point to a circle is always equal to the radius of the circle. Reason (R): Tangents drawn from an external point to a circle are equal in length.

\textbf{Answer:}
(D)

\textbf{Solution:}
\begin{enumerate}
  \item The length of the tangent from an external point is not necessarily equal to the radius.
  \item The theorem states that tangents drawn from an external point to a circle are equal in length.
  \item Therefore, Assertion is false and Reason is true.
\end{enumerate}
\textbf{Derivation:}
\begin{center}
\fitmath{A: \text{False}, R: \text{True}}
\end{center}
\hfill \textit{Circles — Tangents to a circle}
\vspace{0.3cm}

\question \textbf{4.} Assertion (A): If the angle between two tangents drawn from a point P to a circle of radius $r$ and center O is $60^\circ$, then $OP = 2r$. Reason (R): The line segment joining the center to the external point bisects the angle between the two tangents.

\textbf{Answer:}
(A)

\textbf{Solution:}
\begin{enumerate}
  \item In $\triangle OAP$ (where A is point of contact), $\angle OPA = 60^\circ/2 = 30^\circ$.
  \item $\sin(30^\circ) = r/OP$, so $1/2 = r/OP$, which implies $OP = 2r$.
  \item The reason explains the bisection property used in the calculation.
\end{enumerate}
\textbf{Derivation:}
\begin{center}
\fitmath{\begin{aligned}
\sin(30^\circ) = \frac{r}{OP} \\
\implies OP = 2r
\end{aligned}}
\end{center}
\hfill \textit{Circles — Tangents to a circle}
\vspace{0.3cm}

\question \textbf{5.} Assertion (A): A tangent to a circle is perpendicular to the radius at the point of contact. Reason (R): The perpendicular distance from the center of the circle to the tangent is equal to the radius of the circle.

\textbf{Answer:}
(A)

\textbf{Solution:}
\begin{enumerate}
  \item The theorem of tangents states that the tangent at any point of a circle is perpendicular to the radius through the point of contact.
  \item The shortest distance from the center to the line (tangent) is the radius, which confirms the perpendicularity.
  \item Both statements are true and R explains A.
\end{enumerate}
\textbf{Derivation:}
\begin{center}
\fitmath{\text{Radius} \perp \text{Tangent}}
\end{center}
\hfill \textit{Circles — Tangents to a circle}
\vspace{0.3cm}

\question \textbf{6.} Assertion (A): The number of tangents that can be drawn from a point inside a circle is zero. Reason (R): A tangent to a circle intersects the circle at only one point.

\textbf{Answer:}
(A)

\textbf{Solution:}
\begin{enumerate}
  \item A point inside a circle cannot have a line touching the circle at only one point.
  \item A tangent is defined as a line intersecting the circle at exactly one point.
  \item Both are true and R explains why no tangents can exist from an internal point.
\end{enumerate}
\textbf{Derivation:}
\begin{center}
\fitmath{\begin{aligned}
\text{Internal point} \\
\implies 0 \text{ tangents}
\end{aligned}}
\end{center}
\hfill \textit{Circles — Tangents to a circle}
\vspace{0.3cm}

\question \textbf{7.} Assertion (A): If the distance between two parallel tangents of a circle of radius $3$ cm is $6$ cm. Reason (R): Parallel tangents touch the circle at the ends of the diameter.

\textbf{Answer:}
(A)

\textbf{Solution:}
\begin{enumerate}
  \item The distance between two parallel tangents is equal to the diameter of the circle.
  \item Diameter = $2 \times$ radius = $2 \times 3 = 6$ cm.
  \item Both are true and R explains why the distance is the diameter.
\end{enumerate}
\textbf{Derivation:}
\begin{center}
\fitmath{d = 2r = 2(3) = 6}
\end{center}
\hfill \textit{Circles — Tangents to a circle}
\vspace{0.3cm}

\question \textbf{8.} If a tangent $PA$ and $PB$ from a point $P$ to a circle with center $O$ are inclined to each other at an angle of $80^{\circ}$, then find $\angle POA$.

\textbf{Answer:}
\begin{center}
\fitmath{50^{\circ}}
\end{center}

\textbf{Solution:}
\begin{enumerate}
  \item In quadrilateral $\text{OAPB}$, $\angle OAP = \angle OBP = 90^{\circ}$ because the radius is perpendicular to the tangent at the point of contact.
  \item In $\triangle OAP$ and $\triangle OBP$, $OP$ bisects $\angle APB$, so $\angle APO = 40^{\circ}$. In $\triangle OAP$, $\angle POA = 180^{\circ} - (90^{\circ} + 40^{\circ}) = 50^{\circ}$.
\end{enumerate}
\textbf{Derivation:}
\begin{center}
\fitmath{\begin{aligned}
\angle APB = 80^{\circ} \\
\implies \angle APO = 40^{\circ}. \text{ In } \triangle OAP, \angle OAP = 90^{\circ}. \angle POA = 180^{\circ} - 90^{\circ} - 40^{\circ} = 50^{\circ}.
\end{aligned}}
\end{center}
\hfill \textit{Circles — Tangents to a circle}
\vspace{0.3cm}

\question \textbf{9.} A quadrilateral $\text{ABCD}$ is drawn to circumscribe a circle. If $AB = 6\text{ cm}$, $BC = 7\text{ cm}$, and $CD = 4\text{ cm}$, find the length of $AD$.

\textbf{Answer:}
\begin{center}
\fitmath{3\text{ cm}}
\end{center}

\textbf{Solution:}
\begin{enumerate}
  \item For a quadrilateral circumscribing a circle, the sums of opposite sides are equal: $AB + CD = BC + AD$.
  \item Substitute the given values: $6 + 4 = 7 + AD$. Solving for $AD$ gives $10 - 7 = 3$.
\end{enumerate}
\textbf{Derivation:}
\begin{center}
\fitmath{\begin{aligned}
AB + CD = BC + AD \\
\implies 6 + 4 = 7 + AD \\
\implies AD = 10 - 7 = 3 \text{ cm}.
\end{aligned}}
\end{center}
\hfill \textit{Circles — Properties of tangents}
\vspace{0.3cm}

\question \textbf{10.} From a point $Q$, the length of the tangent to a circle is $24\text{ cm}$ and the distance of $Q$ from the center is $25\text{ cm}$. Find the radius of the circle.

\textbf{Answer:}
\begin{center}
\fitmath{7\text{ cm}}
\end{center}

\textbf{Solution:}
\begin{enumerate}
  \item Let $O$ be the center and $A$ be the point of contact. $\triangle OAQ$ is a right-angled triangle with $\angle OAQ = 90^{\circ}$.
  \item By Pythagoras theorem, $OQ^2 = OA^2 + AQ^2$. Substitute $25^2 = r^2 + 24^2$ and solve for $r$.
\end{enumerate}
\textbf{Derivation:}
\begin{center}
\fitmath{\begin{aligned}
25^2 = r^2 + 24^2 \\
\implies 625 = r^2 + 576 \\
\implies r^2 = 49 \\
\implies r = 7 \text{ cm}.
\end{aligned}}
\end{center}
\hfill \textit{Circles — Tangents to a circle}
\vspace{0.3cm}

\question \textbf{11.} Two concentric circles are of radii $5\text{ cm}$ and $3\text{ cm}$. Find the length of the chord of the larger circle which touches the smaller circle.

\textbf{Answer:}
\begin{center}
\fitmath{8\text{ cm}}
\end{center}

\textbf{Solution:}
\begin{enumerate}
  \item Let $AB$ be the chord of the larger circle touching the smaller circle at $M$. $OM \perp AB$ and $OM = 3\text{ cm}$, $OA = 5\text{ cm}$.
  \item In right $\triangle OMA$, $AM = \sqrt{OA^2 - OM^2} = \sqrt{25 - 9} = 4$. Since $M$ is the midpoint, $AB = 2 \times 4 = 8$.
\end{enumerate}
\textbf{Derivation:}
\begin{center}
\fitmath{\begin{aligned}
AM = \sqrt{5^2 - 3^2} = \sqrt{16} = 4 \\
\implies AB = 2 \times 4 = 8 \text{ cm}.
\end{aligned}}
\end{center}
\hfill \textit{Circles — Tangents to a circle}
\vspace{0.3cm}

\question \textbf{12.} If the angle between two tangents drawn from an external point $P$ to a circle of radius $r$ and center $O$ is $60^{\circ}$, find the length of $OP$.

\textbf{Answer:}
\begin{center}
\fitmath{2r}
\end{center}

\textbf{Solution:}
\begin{enumerate}
  \item Let the tangents be $PA$ and $PB$. $OP$ bisects $\angle APB$, so $\angle APO = 30^{\circ}$.
  \item In $\triangle OAP$, $\sin(30^{\circ}) = \frac{OA}{OP} = \frac{r}{OP}$. Since $\sin(30^{\circ}) = 1/2$, we get $OP = 2r$.
\end{enumerate}
\textbf{Derivation:}
\begin{center}
\fitmath{\begin{aligned}
\sin(30^{\circ}) = \frac{r}{OP} \\
\implies \frac{1}{2} = \frac{r}{OP} \\
\implies OP = 2r.
\end{aligned}}
\end{center}
\hfill \textit{Circles — Tangents to a circle}
\vspace{0.3cm}

\question \textbf{13.} Two concentric circles are of radii $5 \text{ cm}$ and $3 \text{ cm}$. Find the length of the chord of the larger circle which touches the smaller circle.

\textbf{Answer:}
\begin{center}
\fitmath{8 \text{ cm}}
\end{center}

\textbf{Solution:}
\begin{enumerate}
  \item Let O be the center, AB be the chord of the larger circle tangent to the smaller circle at point P.
  \item OP is the radius of the smaller circle, so OP = 3 cm and OP $\perp$ AB.
  \item OA is the radius of the larger circle, so OA = 5 cm.
  \item \(
\text{In right triangle OPA}, AP^{2} = OA^{2} - OP^{2} = 5^{2} - 3^{2} = 16, \text{so AP} = 4 cm.
\)
  \item Since the perpendicular from the center bisects the chord, AB = 2 * AP = 8 cm.
\end{enumerate}
\textbf{Derivation:}
\begin{center}
\fitmath{AP = \sqrt{OA^2 - OP^2} = \sqrt{25 - 9} = \sqrt{16} = 4 \text{ cm}; AB = 2 \times 4 = 8 \text{ cm}}
\end{center}
\hfill \textit{Circles — Tangent properties}
\vspace{0.3cm}

\question \textbf{14.} A quadrilateral ABCD is drawn to circumscribe a circle. Prove that $AB + CD = AD + BC$.

\textbf{Answer:}
\begin{center}
\fitmath{AB + CD = AD + BC}
\end{center}

\textbf{Solution:}
\begin{enumerate}
  \item Use the property that lengths of tangents drawn from an external point to a circle are equal.
  \item Let the circle touch sides AB, BC, CD, DA at P, Q, R, S respectively.
  \item \(
AP=AS, BP=BQ, CR=CQ, DR=DS.
\)
  \item Summing the left and right sides of these equations leads to the result.
\end{enumerate}
\textbf{Derivation:}
\begin{center}
\fitmath{AB+CD = (AP+PB)+(CR+RD) = AS+BQ+CQ+DS = (AS+DS)+(BQ+CQ) = AD+BC}
\end{center}
\hfill \textit{Circles — Tangents to a circle}
\vspace{0.3cm}

\question \textbf{15.} If a tangent $PA$ and $PB$ from a point $P$ to a circle with centre $O$ are inclined to each other at an angle of $80^{\circ}$, find $\angle POA$.

\textbf{Answer:}
\begin{center}
\fitmath{50^{\circ}}
\end{center}

\textbf{Solution:}
\begin{enumerate}
  \item In quadrilateral OAPB, $\angle OAP = \angle OBP = 90^{\circ}$ (tangent is perpendicular to radius).
  \item Sum of angles in quadrilateral is $360^{\circ}$, so $\angle AOB = 180^{\circ} - 80^{\circ} = 100^{\circ}$.
  \item In $\triangle OAP$ and $\triangle OBP$, by RHS congruence, $\angle POA = \angle POB$.
  \item Therefore, $\angle POA = \frac{1}{2} \angle AOB = 50^{\circ}$.
\end{enumerate}
\textbf{Derivation:}
\begin{center}
\fitmath{\angle AOB = 180^\circ - 80^\circ = 100^\circ; \angle POA = \frac{100^\circ}{2} = 50^\circ}
\end{center}
\hfill \textit{Circles — Tangent properties}
\vspace{0.3cm}

\question \textbf{16.} Prove that the parallelogram circumscribing a circle is a rhombus.

\textbf{Answer:}
\begin{center}
\fitmath{AB = BC = CD = DA}
\end{center}

\textbf{Solution:}
\begin{enumerate}
  \item Since it is a parallelogram, opposite sides are equal: AB = CD and AD = BC.
  \item From the property of circumscribed quadrilaterals, AB + CD = AD + BC.
  \item Substitute CD with AB and BC with AD: 2AB = 2AD, hence AB = AD.
  \item Since adjacent sides are equal, the parallelogram is a rhombus.
\end{enumerate}
\textbf{Derivation:}
\begin{center}
\fitmath{\begin{aligned}
AB+CD=AD+BC \\
\implies 2AB=2AD \\
\implies AB=AD
\end{aligned}}
\end{center}
\hfill \textit{Circles — Tangents to a circle}
\vspace{0.3cm}

\question \textbf{17.} A point $P$ is $13 \text{ cm}$ from the centre of a circle. The length of the tangent drawn from $P$ to the circle is $12 \text{ cm}$. Find the radius of the circle.

\textbf{Answer:}
\begin{center}
\fitmath{5 \text{ cm}}
\end{center}

\textbf{Solution:}
\begin{enumerate}
  \item Let O be the center and A be the point of contact on the circle.
  \item Triangle OAP is a right-angled triangle at A because the radius is perpendicular to the tangent.
  \item \(
\text{Hypotenuse OP} = 13 cm, \text{tangent AP} = 12 cm.
\)
  \item Using Pythagoras theorem, OA$^{2}$ = OP$^{2}$ - AP$^{2}$ = 13$^{2}$ - 12$^{2}$ = 169 - 144 = 25.
  \item \(
OA = 5 cm.
\)
\end{enumerate}
\textbf{Derivation:}
\begin{center}
\fitmath{OA = \sqrt{13^2 - 12^2} = \sqrt{169 - 144} = \sqrt{25} = 5 \text{ cm}}
\end{center}
\hfill \textit{Circles — Tangents to a circle}
\vspace{0.3cm}

\question \textbf{18.} Prove that the lengths of tangents drawn from an external point to a circle are equal.

\textbf{Answer:}
\begin{center}
\fitmath{PA = PB}
\end{center}

\textbf{Solution:}
\begin{enumerate}
  \item Draw a circle with center O and an external point P.
  \item Draw tangents PA and PB touching the circle at A and B respectively.
  \item Join OA, OB, and OP.
  \item Consider triangles OAP and OBP.
  \item Use the property that the tangent at any point of a circle is perpendicular to the radius through the point of contact.
  \item Use RHS congruence criterion to prove triangles are congruent.
  \item Conclude that corresponding parts of congruent triangles are equal.
\end{enumerate}
\textbf{Derivation:}
In  $\triangle$ OAP  and  $\triangle$ OBP : \textbackslash{}\textbackslash{} 1.  $\angle$ OAP = $\angle$ OBP = 90\textasciicircum{}$\circ  ($Tangent  $\perp$  radius) \textbackslash{}\textbackslash{} 2.  OP = OP  (Common side) \textbackslash{}\textbackslash{} 3.  OA = OB  (Radii of the same circle) \textbackslash{}\textbackslash{} By RHS congruence rule,  $\triangle$ OAP $\cong \triangle$ OBP . \textbackslash{}\textbackslash{} Therefore,  PA = PB  (CPCT).
\hfill \textit{Circles — Tangents to a circle}
\vspace{0.3cm}

\question \textbf{19.} Two concentric circles are of radii $5\text{ cm}$ and $3\text{ cm}$. Find the length of the chord of the larger circle which touches the smaller circle.

\textbf{Answer:}
\begin{center}
\fitmath{8 cm}
\end{center}

\textbf{Solution:}
\begin{enumerate}
  \item Let O be the center of both circles.
  \item Let AB be the chord of the larger circle touching the smaller circle at point P.
  \item OP is the radius of the smaller circle, so $OP = 3\text{ cm}$.
  \item OA is the radius of the larger circle, so $OA = 5\text{ cm}$.
  \item Since AB is a tangent to the inner circle, $OP \perp AB$.
  \item In right $\triangle OPA$, use Pythagoras theorem to find AP.
  \item Since the perpendicular from the center bisects the chord, $AB = 2 \times AP$.
\end{enumerate}
\textbf{Derivation:}
\begin{center}
\fitmath{\begin{aligned}
In \triangle OPA: OA^2 = OP^2 + AP^2 \\
5^2 = 3^2 + AP^2 \\
25 = 9 + AP^2 \\
AP^2 = 16 \\
\Rightarrow AP = 4\text{ cm}. \\
\text{Since OP} \perp AB, AP = PB = 4\text{ cm}. \\
AB = AP + PB = 4 + 4 = 8\text{ cm}.
\end{aligned}}
\end{center}
\hfill \textit{Circles — Properties of tangents}
\vspace{0.3cm}

\question \textbf{20.} A quadrilateral $\text{ABCD}$ is drawn to circumscribe a circle. Prove that $AB + CD = AD + BC$.

\textbf{Answer:}
\begin{center}
\fitmath{AB + CD = AD + BC}
\end{center}

\textbf{Solution:}
\begin{enumerate}
  \item Identify the points of contact on sides AB, BC, CD, and DA as P, Q, R, and S respectively.
  \item Use the theorem that lengths of tangents from an external point to a circle are equal.
  \item Express segments AP, BP, BQ, CQ, CR, DR, DS, and AS in terms of equality.
  \item Sum the equations formed by these segments.
  \item Rearrange the terms to group segments forming the sides of the quadrilateral.
\end{enumerate}
\textbf{Derivation:}
AP = AS (1), BP = BQ (2), CR = CQ (3), DR = DS (4) \textbackslash{}\textbackslash{} Add (1), (2), (3), and (4): \textbackslash{}\textbackslash{} (AP + BP) + (CR + DR) = (AS + DS) + (BQ + CQ) \textbackslash{}\textbackslash{} AB + CD = AD + BC.
\hfill \textit{Circles — Tangents to a circle}
\vspace{0.3cm}

\question \textbf{21.} If a hexagon $\text{ABCDEF}$ circumscribes a circle, prove that $AB + CD + EF = BC + DE + FA$.

\textbf{Answer:}
\begin{center}
\fitmath{AB + CD + EF = BC + DE + FA}
\end{center}

\textbf{Solution:}
\begin{enumerate}
  \item Let the points of contact be P, Q, R, S, T, and U on sides AB, BC, CD, DE, EF, and FA respectively.
  \item By the theorem of tangents from an external point, $AP=AF$, $BP=BQ$, $CQ=CR$, $DR=DS$, $ES=ET$, $FT=FU$ is incorrect; rather, label them appropriately.
  \item Correct logic: $AP=AU, BP=BQ, CQ=CR, DR=DS, ES=ET, FT=FU$.
  \item Sum the left side $AB + CD + EF$ using these relations.
  \item Observe that the sum equals the sum of the remaining sides $BC + DE + FA$.
\end{enumerate}
\textbf{Derivation:}
Let points of contact be  P\_1, P\_2, P\_3, P\_4, P\_5, P\_6 . \textbackslash{}\textbackslash{}  AB+CD+EF = (AP\_1+P\_1B) + (CP\_3+P\_3D) + (EP\_5+P\_5F) = (AP\_6+P\_2B) + (CP\_2+P\_4D) + (EP\_4+P\_6F) \textbackslash{}\textbackslash{} Rearranging: (P\_2B+CP\_2) + (P\_4D+EP\_4) + (P\_6F+AP\_6) = BC + DE + FA .
\hfill \textit{Circles — Tangents to a circle}
\vspace{0.3cm}

\question \textbf{22.} A tower stands vertically on the ground. From a point on the ground, which is $15$ m away from the foot of the tower, the angle of elevation of the top of the tower is found to be $60^\circ$. The height of the tower is:

\textbf{Answer:}
(C)

\textbf{Solution:}
\begin{enumerate}
  \item Let the height of the tower be $h$.
  \item In the right triangle, $\tan(60^\circ) = \frac{\text{height}}{\text{base}}$.
  \item Substitute the values: $\sqrt{3} = \frac{h}{15}$.
  \item Solve for $h$.
\end{enumerate}
\textbf{Derivation:}
\begin{center}
\fitmath{\begin{aligned}
\tan(60^\circ) = \frac{h}{15} \\
\implies \sqrt{3} = \frac{h}{15} \\
\implies h = 15\sqrt{3} \text{ m}
\end{aligned}}
\end{center}
\hfill \textit{Some Applications of Trigonometry — Heights and Distances}
\vspace{0.3cm}

\question \textbf{23.} A kite is flying at a height of $60$ m above the ground. The string attached to the kite is temporarily tied to a point on the ground. The inclination of the string with the ground is $30^\circ$. The length of the string is:

\textbf{Answer:}
(D)

\textbf{Solution:}
\begin{enumerate}
  \item Let $L$ be the length of the string.
  \item In the right triangle, $\sin(30^\circ) = \frac{\text{height}}{\text{hypotenuse}}$.
  \item Substitute the values: $\frac{1}{2} = \frac{60}{L}$.
  \item Solve for $L$.
\end{enumerate}
\textbf{Derivation:}
\begin{center}
\fitmath{\begin{aligned}
\sin(30^\circ) = \frac{60}{L} \\
\implies \frac{1}{2} = \frac{60}{L} \\
\implies L = 120 \text{ m}
\end{aligned}}
\end{center}
\hfill \textit{Some Applications of Trigonometry — Heights and Distances}
\vspace{0.3cm}

\question \textbf{24.} The angle of elevation of the top of a pole from a point on the ground $20$ m away from the foot of the pole is $45^\circ$. The height of the pole is:

\textbf{Answer:}
(B)

\textbf{Solution:}
\begin{enumerate}
  \item Let $h$ be the height of the pole.
  \item Use the relation $\tan(45^\circ) = \frac{h}{20}$.
  \item Since $\tan(45^\circ) = 1$, we have $1 = \frac{h}{20}$.
\end{enumerate}
\textbf{Derivation:}
\begin{center}
\fitmath{\begin{aligned}
\tan(45^\circ) = \frac{h}{20} \\
\implies 1 = \frac{h}{20} \\
\implies h = 20 \text{ m}
\end{aligned}}
\end{center}
\hfill \textit{Some Applications of Trigonometry — Heights and Distances}
\vspace{0.3cm}

\question \textbf{25.} If the length of the shadow of a tower is $\sqrt{3}$ times the height of the tower, then the angle of elevation of the sun is:

\textbf{Answer:}
(A)

\textbf{Solution:}
\begin{enumerate}
  \item Let the height of the tower be $h$. Then the shadow length is $h\sqrt{3}$.
  \item Let $\theta$ be the angle of elevation.
  \item Use $\tan(\theta) = \frac{h}{h\sqrt{3}}$.
  \item Simplify to find $\theta$.
\end{enumerate}
\textbf{Derivation:}
\begin{center}
\fitmath{\begin{aligned}
\tan(\theta) = \frac{h}{h\sqrt{3}} = \frac{1}{\sqrt{3}} \\
\implies \theta = 30^\circ
\end{aligned}}
\end{center}
\hfill \textit{Some Applications of Trigonometry — Heights and Distances}
\vspace{0.3cm}

\question \textbf{26.} An observer $1.5$ m tall is $28.5$ m away from a chimney. The angle of elevation of the top of the chimney from her eyes is $45^\circ$. The height of the chimney is:

\textbf{Answer:}
(C)

\textbf{Solution:}
\begin{enumerate}
  \item Height of observer = $1.5$ m. Distance = $28.5$ m.
  \item Height of chimney above observer's eye level = $28.5 \times \tan(45^\circ) = 28.5$ m.
  \item Total height = $28.5 + 1.5$.
\end{enumerate}
\textbf{Derivation:}
\begin{center}
\fitmath{H = 28.5 + 1.5 = 30 \text{ m}}
\end{center}
\hfill \textit{Some Applications of Trigonometry — Heights and Distances}
\vspace{0.3cm}

\question \textbf{27.} Assertion (A): The angle of elevation of the top of a tower from a point on the ground, which is $30 \text{ m}$ away from the foot of the tower of height $30 \sqrt{3} \text{ m}$, is $60^\circ$. 
Reason (R): In a right-angled triangle, $\tan(\theta) = \frac{\text{Opposite}}{\text{Adjacent}}$.

\textbf{Answer:}
(A)

\textbf{Solution:}
\begin{enumerate}
  \item Identify the height (opposite) = $30\sqrt{3}$ and base (adjacent) = $30$.
  \item Use the formula $\tan(\theta) = \frac{30\sqrt{3}}{30} = \sqrt{3}$.
  \item Since $\tan(60^\circ) = \sqrt{3}$, the angle is $60^\circ$.
\end{enumerate}
\textbf{Derivation:}
\begin{center}
\fitmath{\begin{aligned}
\tan(\theta) = \frac{30\sqrt{3}}{30} = \sqrt{3} \\
\implies \theta = 60^\circ
\end{aligned}}
\end{center}
\hfill \textit{Some Applications of Trigonometry — Angle of Elevation}
\vspace{0.3cm}

\question \textbf{28.} Assertion (A): If the length of the shadow of a vertical pole is equal to its height, then the angle of elevation of the Sun is $45^\circ$. 
Reason (R): The angle of elevation is defined only when the object is taller than the observer.

\textbf{Answer:}
(C)

\textbf{Solution:}
\begin{enumerate}
  \item Let height = $h$ and shadow = $h$.
  \item $\tan(\theta) = \frac{h}{h} = 1$.
  \item $\theta = 45^\circ$. Assertion is true.
  \item Reason is false because angle of elevation is defined by the line of sight relative to the horizontal, regardless of relative height.
\end{enumerate}
\textbf{Derivation:}
\begin{center}
\fitmath{\begin{aligned}
\tan(\theta) = \frac{\text{height}}{\text{shadow}} = \frac{h}{h} = 1 \\
\implies \theta = 45^\circ
\end{aligned}}
\end{center}
\hfill \textit{Some Applications of Trigonometry — Angle of Elevation}
\vspace{0.3cm}

\question \textbf{29.} Assertion (A): The angle of depression of an object from a point $h$ meters above the ground is always equal to the angle of elevation of that point from the object on the ground. 
Reason (R): The angle of elevation and depression are alternate interior angles for parallel horizontal lines.

\textbf{Answer:}
(A)

\textbf{Solution:}
\begin{enumerate}
  \item By the properties of parallel lines, the line of sight creates alternate interior angles.
  \item \(
\text{Thus}, \text{angle of elevation} = \text{angle of depression}.
\)
\end{enumerate}
\textbf{Derivation:}
\begin{center}
\fitmath{\angle \text{elevation} = \angle \text{depression} \text{ (Alternate Interior Angles)}}
\end{center}
\hfill \textit{Some Applications of Trigonometry — Angle of Depression}
\vspace{0.3cm}

\question \textbf{30.} Assertion (A): If a ladder of length $10 \text{ m}$ reaches a wall at a height of $5 \text{ m}$, the angle made by the ladder with the ground is $30^\circ$. 
Reason (R): $\sin(\theta) = \frac{\text{Opposite}}{\text{Hypotenuse}}$.

\textbf{Answer:}
(A)

\textbf{Solution:}
\begin{enumerate}
  \item Opposite = $5$, Hypotenuse = $10$.
  \item $\sin(\theta) = \frac{5}{10} = 0.5$.
  \item $\theta = 30^\circ$.
\end{enumerate}
\textbf{Derivation:}
\begin{center}
\fitmath{\begin{aligned}
\sin(\theta) = \frac{5}{10} = \frac{1}{2} \\
\implies \theta = 30^\circ
\end{aligned}}
\end{center}
\hfill \textit{Some Applications of Trigonometry — Trigonometric Ratios}
\vspace{0.3cm}

\question \textbf{31.} Assertion (A): As an observer moves closer to a tower, the angle of elevation of its top increases. 
Reason (R): $\tan(\theta)$ increases as $\theta$ increases from $0^\circ$ to $90^\circ$.

\textbf{Answer:}
(A)

\textbf{Solution:}
\begin{enumerate}
  \item In $\tan(\theta) = \frac{h}{x}$, as $x$ (distance) decreases, $\tan(\theta)$ increases.
  \item As $\tan(\theta)$ increases, $\theta$ must increase.
\end{enumerate}
\textbf{Derivation:}
\begin{center}
\fitmath{\text{As } x \to 0, \tan(\theta) = \frac{h}{x} \to \infty, \text{ so } \theta \to 90^\circ}
\end{center}
\hfill \textit{Some Applications of Trigonometry — Angle of Elevation}
\vspace{0.3cm}

\question \textbf{32.} A ladder $15$ m long makes an angle of $60^\circ$ with the wall. Find the height of the point where the ladder touches the wall.

\textbf{Answer:}
\begin{center}
\fitmath{7.5 m}
\end{center}

\textbf{Solution:}
\begin{enumerate}
  \item Identify the triangle formed by the ladder (hypotenuse), the wall (adjacent side to 60 degrees), and the ground.
  \item Use the trigonometric ratio $\cos(60^\circ) = \frac{\text{adjacent}}{\text{hypotenuse}}$ to find the height.
\end{enumerate}
\textbf{Derivation:}
\begin{center}
\fitmath{\begin{aligned}
\cos(60^\circ) = \frac{h}{15} \\
\implies \frac{1}{2} = \frac{h}{15} \\
\implies h = 7.5 \text{ m}
\end{aligned}}
\end{center}
\hfill \textit{Some Applications of Trigonometry — Trigonometric Ratios}
\vspace{0.3cm}

\question \textbf{33.} The angle of elevation of the top of a tower from a point on the ground, which is $30$ m away from the foot of the tower, is $30^\circ$. Find the height of the tower.

\textbf{Answer:}
\begin{center}
\fitmath{10\sqrt{3} m}
\end{center}

\textbf{Solution:}
\begin{enumerate}
  \item Use the tangent ratio $\tan(\theta) = \frac{\text{opposite}}{\text{adjacent}}$ where $\theta = 30^\circ$.
  \item Substitute the given values and solve for the height $h$.
\end{enumerate}
\textbf{Derivation:}
\begin{center}
\fitmath{\begin{aligned}
\tan(30^\circ) = \frac{h}{30} \\
\implies \frac{1}{\sqrt{3}} = \frac{h}{30} \\
\implies h = \frac{30}{\sqrt{3}} = 10\sqrt{3} \text{ m}
\end{aligned}}
\end{center}
\hfill \textit{Some Applications of Trigonometry — Angle of Elevation}
\vspace{0.3cm}

\question \textbf{34.} A kite is flying at a height of $60$ m above the ground. The string attached to the kite is temporarily tied to a point on the ground. The inclination of the string with the ground is $60^\circ$. Find the length of the string.

\textbf{Answer:}
\begin{center}
\fitmath{40\sqrt{3} m}
\end{center}

\textbf{Solution:}
\begin{enumerate}
  \item Represent the situation as a right triangle where the height is the opposite side and string length is the hypotenuse.
  \item Use $\sin(60^\circ) = \frac{\text{height}}{\text{length}}$ to solve for the length.
\end{enumerate}
\textbf{Derivation:}
\begin{center}
\fitmath{\begin{aligned}
\sin(60^\circ) = \frac{60}{L} \\
\implies \frac{\sqrt{3}}{2} = \frac{60}{L} \\
\implies L = \frac{120}{\sqrt{3}} = 40\sqrt{3} \text{ m}
\end{aligned}}
\end{center}
\hfill \textit{Some Applications of Trigonometry — Real-world problems}
\vspace{0.3cm}

\question \textbf{35.} A $1.5$ m tall boy is standing at some distance from a $30$ m tall building. The angle of elevation from his eyes to the top of the building increases from $30^\circ$ to $60^\circ$ as he walks towards the building. Find the distance he walked towards the building.

\textbf{Answer:}
\begin{center}
\fitmath{19\sqrt{3} m}
\end{center}

\textbf{Solution:}
\begin{enumerate}
  \item Determine the effective height of the building relative to the boy's eye level: $30 - 1.5 = 28.5$ m.
  \item Use $\cot(30^\circ)$ and $\cot(60^\circ)$ to find the difference in ground distances.
\end{enumerate}
\textbf{Derivation:}
\begin{center}
\fitmath{d = 28.5(\cot(30^\circ) - \cot(60^\circ)) = 28.5(\sqrt{3} - \frac{1}{\sqrt{3}}) = 28.5(\frac{2}{\sqrt{3}}) = \frac{57}{\sqrt{3}} = 19\sqrt{3} \text{ m}}
\end{center}
\hfill \textit{Some Applications of Trigonometry — Angle of Elevation}
\vspace{0.3cm}

\question \textbf{36.} From a point on the ground, the angles of elevation of the bottom and top of a transmission tower fixed at the top of a $20$ m high building are $45^\circ$ and $60^\circ$ respectively. Find the height of the tower.

\textbf{Answer:}
\begin{center}
\fitmath{20(\sqrt{3}-1) m}
\end{center}

\textbf{Solution:}
\begin{enumerate}
  \item Let $h$ be the height of the tower. The building height is $20$.
  \item Use $\tan(45^\circ)$ to find the distance from the point to the building, then use $\tan(60^\circ)$ for the total height.
\end{enumerate}
\textbf{Derivation:}
\begin{center}
\fitmath{\begin{aligned}
\text{Distance } x = \frac{20}{\tan(45^\circ)} = 20. \text{ Total height } (20+h) = x \tan(60^\circ) = 20\sqrt{3}. \\
\implies h = 20\sqrt{3} - 20 = 20(\sqrt{3}-1) \text{ m}
\end{aligned}}
\end{center}
\hfill \textit{Some Applications of Trigonometry — Multiple Angles of Elevation}
\vspace{0.3cm}

\question \textbf{37.} A tower stands vertically on the ground. From a point on the ground, which is $15$ m away from the foot of the tower, the angle of elevation of the top of the tower is found to be $60^\circ$. Find the height of the tower.

\textbf{Answer:}
\begin{center}
\fitmath{15\sqrt{3} \text{ m}}
\end{center}

\textbf{Solution:}
\begin{enumerate}
  \item Let $h$ be the height of the tower and $d = 15$ m be the distance from the foot.
  \item Use the tangent ratio: $\tan(\theta) = \frac{\text{Opposite}}{\text{Adjacent}}$.
  \item Substitute the values: $\tan(60^\circ) = \frac{h}{15}$.
  \item Solve for $h$: $h = 15 \times \sqrt{3}$.
\end{enumerate}
\textbf{Derivation:}
\begin{center}
\fitmath{\begin{aligned}
\tan(60^\circ) = \frac{h}{15} \\
\implies \sqrt{3} = \frac{h}{15} \\
\implies h = 15\sqrt{3}
\end{aligned}}
\end{center}
\hfill \textit{Some Applications of Trigonometry — Heights and Distances}
\vspace{0.3cm}

\question \textbf{38.} A $1.5$ m tall boy is standing at some distance from a $30$ m tall building. The angle of elevation from his eyes to the top of the building increases from $30^\circ$ to $60^\circ$ as he walks towards the building. Find the distance he walked towards the building.

\textbf{Answer:}
\begin{center}
\fitmath{19\sqrt{3} \text{ m}}
\end{center}

\textbf{Solution:}
\begin{enumerate}
  \item Height of building above boy's eyes = $30 - 1.5 = 28.5$ m.
  \item Let $x_1$ be the distance at $30^\circ$ and $x_2$ be the distance at $60^\circ$.
  \item Calculate $x_1 = 28.5 \cot(30^\circ)$ and $x_2 = 28.5 \cot(60^\circ)$.
  \item Distance walked = $x_1 - x_2 = 28.5(\sqrt{3} - \frac{1}{\sqrt{3}}) = 28.5(\frac{2}{\sqrt{3}}) = 19\sqrt{3}$.
\end{enumerate}
\textbf{Derivation:}
\begin{center}
\fitmath{\begin{aligned}
x_1 = 28.5\sqrt{3}, x_2 = \frac{28.5}{\sqrt{3}} \\
\implies \text{dist} = 28.5(\frac{3-1}{\sqrt{3}}) = \frac{57}{\sqrt{3}} = 19\sqrt{3}
\end{aligned}}
\end{center}
\hfill \textit{Some Applications of Trigonometry — Two angles of elevation}
\vspace{0.3cm}

\question \textbf{39.} From the top of a $7$ m high building, the angle of elevation of the top of a cable tower is $60^\circ$ and the angle of depression of its foot is $45^\circ$. Determine the height of the tower.

\textbf{Answer:}
\begin{center}
\fitmath{7(1+\sqrt{3}) \text{ m}}
\end{center}

\textbf{Solution:}
\begin{enumerate}
  \item Let the building height $AB = 7$ m. The tower height $CD = h$.
  \item From the triangle at the base, the distance between building and tower is $7 \cot(45^\circ) = 7$ m.
  \item In the upper triangle, the height above the building level is $h' = 7 \tan(60^\circ) = 7\sqrt{3}$.
  \item Total height $h = 7 + 7\sqrt{3}$.
\end{enumerate}
\textbf{Derivation:}
\begin{center}
\fitmath{h = 7 + 7\tan(60^\circ) = 7(1 + \sqrt{3})}
\end{center}
\hfill \textit{Some Applications of Trigonometry — Angle of elevation and depression}
\vspace{0.3cm}

\question \textbf{40.} A kite is flying at a height of $60$ m above the ground. The string attached to the kite is temporarily tied to a point on the ground. The inclination of the string with the ground is $60^\circ$. Find the length of the string, assuming that there is no slack in the string.

\textbf{Answer:}
\begin{center}
\fitmath{40\sqrt{3} \text{ m}}
\end{center}

\textbf{Solution:}
\begin{enumerate}
  \item Height $h = 60$ m, angle $\theta = 60^\circ$. Let string length be $L$.
  \item Use sine ratio: $\sin(60^\circ) = \frac{\text{Opposite}}{\text{Hypotenuse}} = \frac{60}{L}$.
  \item Substitute $\sin(60^\circ) = \frac{\sqrt{3}}{2}$.
  \item Solve: $L = \frac{60 \times 2}{\sqrt{3}} = \frac{120}{\sqrt{3}} = 40\sqrt{3}$.
\end{enumerate}
\textbf{Derivation:}
\begin{center}
\fitmath{\begin{aligned}
\sin(60^\circ) = \frac{60}{L} \\
\implies \frac{\sqrt{3}}{2} = \frac{60}{L} \\
\implies L = \frac{120}{\sqrt{3}} = 40\sqrt{3}
\end{aligned}}
\end{center}
\hfill \textit{Some Applications of Trigonometry — Sine ratio application}
\vspace{0.3cm}

\end{questions}

\end{document}
